% !TeX program = xelatex
\documentclass{article}
\usepackage{kotex}
\usepackage[a4paper, left=3cm, right=3cm, top=2.5cm, bottom=2.5cm]{geometry}
\usepackage{amsmath, amssymb}
\usepackage{tcolorbox}
\usepackage{caption}
\usepackage{setspace}
\usepackage{xcolor}
\usepackage{titlesec}
\usepackage{tikz}
\usetikzlibrary{arrows.meta, calc}

\setstretch{1.3}
\renewcommand{\figurename}{그림}
\titleformat{name=\section,numberless}
  {\normalfont\large\bfseries\color{blue}}{}
  {0pt}{}

\title{2.1 측지선}
\date{}
\author{}

\begin{document}
\maketitle

\begin{tcolorbox}[colback=blue!5, colframe=blue!60!black,
  title=제 2 장\quad 구면 기하학 --- 한눈에 보기,
  boxrule=0.5mm, arc=3mm]
\begin{tabular}{@{}ll@{}}
2.1  & 측지선 \\
2.2  & 구면의 측지선 \\
2.3  & 구면 삼각형의 여섯 개의 각 \\
2.4  & 변에 대한 코사인 법칙 \\
2.5  & 쌍대 구면 삼각형 \\
2.6  & 각에 대한 코사인 법칙 \\
2.7  & 구면 삼각형에 대한 사인 법칙 \\
2.10 & 입체 사영의 응용 \\
\end{tabular}
\end{tcolorbox}

\section*{2.1 측지선}

\begin{tcolorbox}[
  colback=teal!4, colframe=teal!65!black,
  title=정의 \& 핵심 규칙,
  fonttitle=\bfseries, boxrule=0.5mm, arc=3mm]

\textbf{\textcolor{teal!65!black}{정의 2.1.1 (측지선).}}\quad
공간에서 두 점을 잇는 \textbf{가장 짧은 곡선}은 그 공간에서의 \textbf{측지선(geodesic)}이다.

\noindent\textcolor{teal!50!black}{\rule{\linewidth}{0.4pt}}

\textbf{\textcolor{teal!65!black}{핵심 규칙.}}\quad
측지선을 찾을 때, 공간 $S$에 \textbf{온전히 놓여 있는 곡선만} 고려한다.
$S$를 벗어나는 더 짧은 곡선이 있더라도 측지선과 무관하다 --- $S$를 마치 우주 전체라 생각한다.

\noindent\textcolor{teal!50!black}{\rule{\linewidth}{0.4pt}}

\textbf{\textcolor{teal!65!black}{정의 2.1.2 (측지 삼각형).}}\quad
세 꼭짓점을 \textbf{측지선}으로 연결하여 만든 ``삼각형''을 \textbf{측지 삼각형}이라 한다.

\end{tcolorbox}

\medskip
\textbf{보기 2.1.1} (지구의 측지선).\quad
지구에 있는 두 점을 잇는 측지선은 두 점 사이에 줄을 연결하고 팽팽하게 잡아당겨서 얻을 수 있다.
로스앤젤레스와 런던을 잇는 측지선은 캐나다의 중앙을 지나 북동쪽으로 진행하다가 그린란드의 남쪽 꼭대기를 통과하면서 동쪽으로 방향을 바꾸어 남동쪽의 방향으로 런던에 도착한다.
배나 항공기는 먼 거리를 갈 때에는 이러한 ``대원 항로''를 따라 운행하여 연료를 절약한다(그림 2.1).

\begin{figure}[h]
\centering
\begin{tikzpicture}[scale=1.1]
  % Sphere outline
  \draw[thick] (0,0) circle (2cm);
  % Equator (ellipse, perspective)
  \draw[thick, dashed] (0,0) ellipse (2cm and 0.45cm);
  % Meridian lines (several great circles as rotated ellipses)
  \foreach \ang in {0, 30, 60, 150}{
    \begin{scope}[rotate=\ang]
      \draw[thin, gray!60] (0,-2) arc (-90:90:0.25cm and 2cm);
    \end{scope}
  }
  % Great-circle geodesic (LA to London): thick tilted ellipse arc
  \begin{scope}[rotate=25]
    \draw[very thick] (-1.8, 0.87) arc (154:26:2cm and 0.8cm);
  \end{scope}
  % City points
  \fill[black] (-0.55,-1.18) circle (2.5pt) node[below left, font=\small]{LA};
  \fill[black] (0.55, 1.22) circle (2.5pt) node[above right, font=\small]{런던};
  % Label
  \node[right, font=\small] at (2.1, 0.3) {측지선};
\end{tikzpicture}
\caption{로스앤젤레스와 런던을 잇는 측지선 (대원호)}
\label{fig:2.1}
\end{figure}

\textbf{보기 2.1.2} (원기둥의 측지선).\quad
종이를 말아서 원기둥을 만들고 두 점을 줄로 연결하고 팽팽하게 잡아당기면 측지선을 얻는다.
줄로 원기둥을 몇 번 감느냐에 따라 무수히 많은 측지선을 얻을 수 있다.
그러나 두 점 사이에서 길이를 최소로 하는 측지선은 유일하다.
이 측지선은 종이를 펼쳤을 때 선분이 된다.

\begin{figure}[h]
\centering
\begin{tikzpicture}[scale=0.9]
  %--- Left: cylinder with helical geodesic ---
  % Cylinder body
  \draw[thick] (-1,0) -- (-1,-2.5);
  \draw[thick] ( 1,0) -- ( 1,-2.5);
  % Top ellipse
  \draw[thick] (0,0) ellipse (1cm and 0.3cm);
  % Bottom ellipse
  \draw[thick] (0,-2.5) ellipse (1cm and 0.3cm);
  % Helical geodesic (sinusoidal curve on the surface)
  \draw[very thick] plot[domain=0:1, samples=60]
    ({sin(360*\x)}, {-0.3*cos(360*\x) - 2.5*\x});
  % Labels
  \node[below, font=\small] at (0,-2.9) {원기둥};

  %--- Right: cone with geodesic ---
  \begin{scope}[xshift=4cm]
    % Cone outline
    \draw[thick] (-1.2,-2.5) -- (0,0) -- (1.2,-2.5);
    % Bottom ellipse of cone
    \draw[thick] (0,-2.5) ellipse (1.2cm and 0.32cm);
    % Geodesic on cone (straight when unrolled)
    \draw[very thick] (-0.9,-1.8) .. controls (-0.1,-0.5) .. (0.85,-1.6);
    % Labels
    \node[below, font=\small] at (0,-2.9) {원뿔};
  \end{scope}
\end{tikzpicture}
\caption{원기둥과 원뿔에서의 측지선}
\label{fig:2.2}
\end{figure}

\medskip
\textbf{보기 2.1.3} (원뿔에서의 측지 삼각형).\quad
중심각의 크기 $\theta$인 부채꼴로부터 얻은 원뿔에서
꼭짓점 $A, B, C$를 갖는 측지 삼각형이 부채꼴 꼭짓점의 내부에 있으면:
\[
  \angle A + \angle B + \angle C = 540° - \theta \tag{2.1}
\]
특히 $\theta < 360°$이면 $\angle A + \angle B + \angle C > 180°$임을 알 수 있다.

\begin{tcolorbox}[
  colback=teal!4, colframe=teal!65!black,
  title=핵심 결론: 내각의 합과 곡면의 모양,
  fonttitle=\bfseries, boxrule=0.5mm, arc=3mm]

\begin{center}
\renewcommand{\arraystretch}{1.3}
\begin{tabular}{|l|l|l|}
\hline
곡면의 모양 & 내각의 합 & 예 \\
\hline
평면 (flat) & $= 180°$ & 유클리드 평면 \\
볼록 (sphere-like) & $> 180°$ & 구면, 원뿔($\theta < 360°$) \\
오목 (saddle-like) & $< 180°$ & 말안장 곡면 \\
\hline
\end{tabular}
\end{center}

이 사실은 이차원 생명체가 자신이 살고 있는 곡면의 모양을
곡면을 벗어나지 않고 \textbf{측지 삼각형의 내각의 합}으로 판단할 수 있음을 보여준다.
유명한 \textbf{가우스-보네 정리}가 이 추측이 참임을 말해준다.

\end{tcolorbox}

\subsection*{연습문제 2.1}

\begin{enumerate}
  \item \textbf{(문제 2.1.1)}\quad
    중심각의 크기가 $0° < \theta < 540°$인 임의의 부채꼴로부터 얻은 곡면에 대해
    식 (2.1)을 증명하여라.
    $\theta \geq 540°$인 경우 곡면에 있는 측지 삼각형의 내부에 부채꼴의 꼭짓점이
    존재할 수 없음을 보여라.
    부채꼴의 꼭짓점이 내부에 없는 경우 측지 삼각형의 내각의 합은 얼마인가?
\end{enumerate}

\end{document}
